% =============================================================================
\section{Python bindings Tool Comparison}
\label{sec:comp}
% =============================================================================
There are several different tools and options for generating Python bindings.  The following is by no means an exhaustive list of such tools; however, it gives a taste of the rich possibilities one faces when there is a need to make C++ code accessible to Python developers. One goal that all these tools share is transforming Python/C API into a relatively more user-friendly interface. Most listed tools are designed to wrap C++ interfaces non-intrusively; that is, the C++ code of an entity does not change when it is wrapped.

% -----------------------------------------------------------------------------
\subsection{Binding Tools}
\label{ssec:comp:tools}
% -----------------------------------------------------------------------------
Several tools that facilitate binding generation have been introduced throughout the years. We have concentrated on the three most propitious tools, namely SWIG, PyBind11, and Boost.Python. (A short description of additional tools appear in Section~\ref{sec:tools}.) We have conducted several experiments to compare the performance of the bindings generated using the selected tools; see Section~\ref{ssec:comp:exp}.

% SWIG
SWIG is a software development tool used to create bindings to C and C++ code for several different target languages, including common scripting languages and compiled languages.\footnote{See \url{http://www.swig.org/}.}
% SWIG is a software development tool that connects programs written in C and C++ with a variety of high-level programming languages.\footnote{See \url{http://www.swig.org/}.}
% It is different than any of the others tools listed here.
% SWIG is used to create bindings to C and C++ code for different target languages, and not only Python, including common scripting languages such as Javascript, Perl, PHP, Python, Tcl, or Ruby. The list of supported languages also includes non-scripting languages such as C\#, D, Go language, Java including Android, Lua, OCaml, Octave, Scilab and R. Also several interpreted and compiled Scheme implementations (Guile, MzScheme/Racket) are supported.
%
SWIG is essentially a compiler that parses external C declarations and creates the wrappers needed to access those declarations. It is most commonly used to create high-level interpreted or compiled programming environments, user interfaces, and as a tool for testing and prototyping C/C++ software. SWIG is typically used to parse C/C++ interfaces and generate the `glue code' required for the above target languages to call into the C/C++ code.
% SWIG can also export its parse tree in the form of XML.
% SWIG is free software and the code that SWIG generates is compatible with both commercial and non-commercial projects.

% Boost.Python
Boost is a large and complex suite of utility libraries that works with almost every C++ compiler in existence. Boost.Python~\cite{bhsbp-ag-03} is a member of the Boost suite that enables seamless interoperability between C++ and the Python programming language.\footnote{See   \url{https://www.boost.org/doc/libs/1_70_0/libs/python/doc/html/index.html}.} Similar to \pyLstinline{cppyy}, Boost.Python focuses at C++; unlike \pyLstinline{cppyy} it uses C++ to specify and build bindings, taking advantage of the metaprogramming tools and polymorphism in C++. The library is rich and can create wrappers for all kind of types and features found in modern C++ code, such as references, pointers, iterators, overloaded functions, and default arguments.

% The library includes support for:
% \begin{itemize}
% \item References and Pointers
% \item Globally Registered Type Coercions
% \item Automatic Cross-Module Type Conversions
% \item Efficient Function Overloading
% \item C++ to Python Exception Translation
% \item Default Arguments
% \item Keyword Arguments
% \item Manipulating Python objects in C++
% \item Exporting C++ Iterators as Python Iterators
% \item Documentation Strings
% \end{itemize}

% PyBind11
PyBind11 is a lightweight header-only library that exposes C++ types in Python and vice versa, mainly to create Python bindings of existing C++ code.\footnote{See \url{https://github.com/pybind/pybind11}.} Its goal, like the goal of Boost.Python, is to minimize boilerplate code in traditional extension modules by inferring type information using compile-time introspection. Its syntax is also similar to the syntax of Boost.Python. In fact, it is a tiny self-contained version of Boost.Python with everything stripped away that is not relevant for binding generation.
% EFI, YOU MAY WISH TO PASS THE FOLLOWING CUT TEXT TO TEH APPENDIX
%Without comments, the core header files only require approximately 4K lines of code and depend on Python and the C++ standard library. This compact implementation was possible thanks to some of the new C++11 language features (specifically, tuples, lambda functions and variadic templates). Since its creation, this library has grown beyond Boost.Python in several ways, leading to simpler binding code in many common situations.

% -----------------------------------------------------------------------------
\subsection{Experiments}
\label{ssec:comp:exp}
% -----------------------------------------------------------------------------
C++ Python bindings for a certain subset of \cgal{} packages are available as part of an experimental project that uses SWIG. While the project has been conceived more than a decade ago, it contains bindings only for a limited subset of \cgal{} packages. Naturally, embarking on the existing venture might have saved time in developing new bindings, on the one hand, and in exploiting the new bindings by our users on the other. Another advantage of using SWIG, and probably the reason for choosing SWIG in the first place in the aforementioned project, is the support for multiple languages. Once bindings have been developed for one language, generating bindings for additional languages does not require additional development resources. However, as our experiments show, this profiteering has its cost---the generated bindings are not as time or space efficient as their counterparts generated with other tools. The PyBind11 tool is relatively new. Both the PyBind11 and the Boost.Python tools are widely spread and tested, and actively maintained and enhanced. They are developed using advanced features of C++11 and higher versions. These two tools share the same interface; that is, the C++ code used to generate the bindings using both tools is similar and in some cases even identical. Switching between using bindings generated by PyBind11 and bindings generated by Boost.Python is seamless---it amounts to replacing one library by another.
%As we are very much familiar with developing code that exploits advanced feature of C++, developing bindings using these tools does not present a major hurdle for us. On the other hand, developing bindings using SWIG, requires writing code that describe the bindings in a proprietary language.  Some developers may avoid using Boost.Python, because it requires dealing with a large and complex system. This does not present a drawback for us, because \cgal{} depends on several other libraries of Boost; thus, we are familiar with Boost and must deal with it anyway.
The compatibility of Boost with almost every C++ compiler in general has its cost; arcane template tricks and workarounds are necessary to support the oldest and buggiest of compiler specimens.  Now that C++11-compatible compilers are widely available, this heavy machinery has become an excessively large dependency. However, our main concern was the efficiency of the generated bindings in terms of execution time and memory space.

We have conducted two sets of experiments. In the first set we compared the time it took to refer to C++ objects from Python code and in the second set we compared the time it took to call C++ functions from Python code. The first set was further divided into two subsets. In the first subset we compared the time it took to make many references to many objects and in the second subset we compared the time it took to make many references to a single object; see Table~\ref{tab:res:refs} for the results.
%
The second set was divided into two subsets as well. In the first subset we compared the time it took to call many small functions, and in the second subset we compared the time it took to call few large functions. We also used the second set of experiments to measure the time it took to compute the various tasks in native Python; see Table~\ref{tab:res:calls} for the results.
%
The results we obtained left nothing in doubt.
% We abandoned using SWIG due do its high overhead and we picked Boost.python for its high efficiency in every category.
We picked Boost.python for its high efficiency in every category.

\begin{table}
  \caption{Time consumption of references measured in seconds.}
  \label{tab:res:refs}
  \centering\begin{tabular}{|l|r|r|r|}
    \hline
    & \textbf{pybind11} & \textbf{boost} & \textbf{SWIG}\\
    \hline
    \hline
    Many references to many objects    & 6.031  & 4.453 & 15.687\\
    \hline
    Many references to a single object & 4.906  & 2.796 &  9.640\\
    \hline
  \end{tabular}
\end{table}
\begin{table}
  \caption{Time consumption of function calls measured in seconds.}
  \label{tab:res:calls}
  \centering\begin{tabular}{|l|r|r|r|r|}
    \hline
    & \textbf{pybind11} & \textbf{boost} & \textbf{SWIG} & \textbf{Native}\\
    \hline
    \hline
    Many small tasks & 26.718 & 13.656 & 13.140 & 13.359\\
    \hline
    One large task   & 0.718  & 0.718  & 0.734  & 17.843\\
    \hline
  \end{tabular}
\end{table}

% Boost.Python handles shared\_ptr and other smart pointers extremely gracefully
